\scanpage{201}
\mspage{117}
% archive sheet number: 200
% element: text | top, lines 1-3 | où M pi désigne
où $M_\pi$ désigne la sous-catégorie str. pleine de $M$ formée des objets $\pi$-accessibles.
Ceci posé, on a les implications
% element: diagram | upper third, line 4 | 1 implies 2 implies 3 implies 4 implies 5
\diagram{\begin{tikzcd}[ampersand replacement=\&]
1)\arrow[r,Rightarrow]\&2)\arrow[r,Rightarrow]\&3)\arrow[r,Rightarrow]\&4)\arrow[r,Rightarrow]\&5),
\end{tikzcd}}
% element: text | upper third, lines 5-8 | et dans le cas particulier
et dans le cas particulier où $M=(\mathrm{Ens})$, $C=$ sous-catégorie pleine formée des ensembles de cardinal $\leq\pi_0$, on a aussi $5)\Rightarrow1)$ i.e. toutes ces conditions sont équivalentes.

% element: text | middle, line 9 | Dém 1 implies 2
\uline{Dém.}\quad $1)\Rightarrow2)$, cf dém. de 4.5.4.

% element: text | middle, lines 10-18 | 2 implies 3 La catégorie
$2)\Rightarrow3)$ La catégorie $J=\mathrm{Filt}^{\pi}_{/F}$ est stable par $\varinjlim$ de cardinal $\leq\pi$, $\mathrm{Filt}^{\pi}$ l'étant, donc $J$ est \unsure{grande|grand} devant $\pi$, d'ailleurs $F=\varinjlim_{j\in J}X_j$, \unsure{dont|où} $j\longmapsto X_j$ est le foncteur source $\mathrm{Filt}^{\pi}_{/F}\longrightarrow\mathrm{Filt}^{\pi}$. \unsure{D'où|D'ou} la repr. voulue de $F$, quitte à remplacer $J$ par un ens ordonné cofinal $\unsure{I|\bar I}$, p.ex. en prenant pour $I$ l'ens des sous-catégories \deleted{\unsure{filtrantes de}} \deleted{\unsure{filtrantes}} $J$ de $I$ \deleted{de cardinal} \deleted{grandes devant $\pi$} \deleted{$\leq\pi$} telles que $\operatorname{card}\mathrm{Fl}\,J\leq\pi$. La majoration $\operatorname{card}I\leq\pi'{}^\pi$, pour $F\in\mathrm{Ob}\,\mathrm{Filt}^{\pi'}$, se fait par AQT.

% element: text | lower third, line 19 | L'implication 3 to 4
L'implication $3)\longrightarrow4)$ est tautologique

% element: text | foot, lines 20-23 | 4 implies 5
$4)\Rightarrow5)$, car si $F\in\mathrm{Ob}\,M_\pi$, on \unsure{trouve|tire} par 4) que $F$ est facteur direct d'un objet de $\mathrm{Filt}^{\pi}$, et $\mathrm{Filt}^{\pi}$ étant stable par $\varinjlim$
% slanted
% element: margin note | bottom left, beside 4 implies 5 | Comparer le lemme
\margin{Comparer le lemme 4.5.14 \unsure{page|pg} 133}
